\section{Marco teórico y trabajos relacionados}

La arquitectura MVC ha sido extensamente estudiada en la literatura de ingeniería de software. Pop y Altar \cite{pop2014designing} presentan un modelo para desarrollo rápido de aplicaciones web basado en MVC, implementado en PHP. Los autores enfatizan la separación de responsabilidades como clave para mejorar el tiempo de desarrollo y mantenimiento, con beneficios en la reutilización de código y desarrollo paralelo.

Aniche et al. \cite{aniche2016satt} proponen métricas y umbrales específicos para aplicaciones web MVC, adaptando las métricas de Chidamber y Kemerer al contexto arquitectural. Su estudio establece umbrales para cohesion (LCOM), acoplamiento (CBO), complejidad (WMC) y otras métricas, diferenciando entre capas Controller, Service, Entity y Repository. Este trabajo es fundamental para evaluar la calidad de código MVC de manera contextualizada.

En un estudio más reciente, Aniche et al. \cite{aniche2018code} identifican seis olores de código específicos para aplicaciones MVC: Promiscuous Controller, Brain Controller, Meddling Service, Brain Repository, Laborious Repository Method y Fat Repository. Estos olores capturan violaciones a los principios de diseño en el contexto MVC, permitiendo detectar problemas de mantenibilidad y escalabilidad.

Necula \cite{necula2024exploring} realiza un análisis amplio de las aplicaciones de mercado y tecnológicas de MVC. La autora destaca la adopción de MVC en frameworks como Spring MVC, ASP.NET MVC y JavaScript (Angular, React, Vue.js), así como su presencia en sectores como e-commerce, servicios financieros y educación. El estudio identifica tendencias emergentes como la transición hacia MVVM y la integración con IA y machine learning.


Otros trabajos relevantes incluyen el uso de MVC en aplicaciones móviles \cite{narsoo2009application} y su aplicación en sistemas empresariales \cite{zhao2010design}. Estos estudios confirman la versatilidad de MVC pero también destacan la necesidad de adaptaciones específicas según el dominio de aplicación.

Nuestro trabajo se diferencia al proporcionar una evaluación empírica completa de un sistema MVC real, combinando métricas de calidad, detección de olores y análisis de arquitectura, contribuyendo así a la validación práctica de los principios teóricos establecidos en la literatura.